% !TeX spellcheck = it_IT

\documentclass[format=169, handout, 10pt]{beamer}

\usepackage[utf8]{inputenc}
\usepackage[italian]{babel}
\usepackage{hyperref}
\usetheme{metropolis}

\usepackage{mathtools}
\DeclarePairedDelimiter{\ceil}{\lceil}{\rceil}

\title{Esercitazioni di Informatica B}
\subtitle{Funzioni in MATLAB}

\author{Stefano Cereda\\
	stefano.cereda@polimi.it
}
\date{/11/2018}
\institute[PoliMi]{\vspace{0.5cm}\centering Politecnico di Milano \\ \vspace{0.2cm}
	\includegraphics[width=0.2\textwidth]{../logopolimi}}

\setbeamercovered{invisible}


\makeindex


\begin{document}
	\begin{frame}
	\maketitle
\end{frame}

\begin{frame}{Funzione split}
Scrivere una funzione (split) che data una variabile di ingresso $m$ (scalare, vettore o
matrice di numeri) ed un numero $x$ restituisce due vettori $vMin$ e $vMax$. $vMin$
contiene tutti gli elementi di $m$ minori di $x$, mentre $vMax$ gli elementi maggiori
di $x$.
\end{frame}

\begin{frame}[allowframebreaks]{Struct e funzioni: codice IBAN}
Un codice IBAN è una sequenza di 27 caratteri alfanumerici così composta:
\begin{itemize}
\item 2 caratteri maiuscoli rappresentanti la sigla internazionale della Nazione
\item 2 cifre (CIN Europeo)
\item 1 carattere maiuscolo (CIN Italiano)
\item 5 cifre (ABI)
\item 5 cifre (CAB)
\item 12 cifre (Numero di Conto)
\end{itemize}
\begin{tabular}{|c|c|c|c|c|c|}
	\hline
	naz & cin\_e & cin\_I & abi & cab	&	conto \\
	\hline
	IT&02&L&12345&12345&123456789012\\
	\hline
\end{tabular}


Si scrivano le seguenti funzioni:
\begin{itemize}
\item \emph{removeSpaces}, che prenda in ingresso una stringa strIn e restituisca una stringa strOut contenente tutti i caratteri di strIn ad esclusione degli spazi.
\item \emph{isAlpha}, che prenda in ingresso una stringa str e restituisca 1 quando str contiene esclusivamente lettere maiuscole, o 0 altrimenti.
\item \emph{isDigit}, che prenda in ingresso una stringa str e restituisca 1 quando str contiene esclusivamente caratteri corrispondenti a cifre, o 0 altrimenti.
\end{itemize}

Si utilizzino tali funzioni per scrivere un’ulteriore funzione \emph{checkIBAN} che riceve come parametro una stringa e restituisce 1 se la stringa inserita, una volta rimossi gli spazi, corrisponde ad un codice IBAN che rispetta lo schema sopra illustrato, o 0 altrimenti.
\end{frame}

	
\begin{frame}[allowframebreaks]{Struct e funzioni: biblioteca}
Si vuole automatizzare la gestione di una biblioteca di un piccolo paese.

Come primo passo si vuole gestire una matrice di dimensioni $N \times 2$ contenente nella
prima colonna un identificativo numerico e nella seconda colonna un numero di
tessera.

Immaginate l’identificativo numerico come un codice interno alla biblioteca mentre il
numero di tessera come il numero di una tessera esistente appartenente (ad esempio
il vostro codice persona del politecnico).

\begin{tabular}{|c|c|}
	\hline
	ID	&	Codice Persona \\
	\hline
	1	&	123456789	\\
	\hline
	2	&	987654321	\\
	\hline
	...	&	... \\
	\hline
\end{tabular} 


Si forniscano le funzioni necessarie all'inserimento di una nuova tessera e alla ricerca
dell'identificativo numerico dato un numero di tessera (l’identificativo numerico è un
numero crescente che parte da 1).

Funzione \emph{inserisci}:\\
\emph{input}: matrice Soci e nuovo numero di tessera\\
\emph{output}: la matrice Soci

Funzione \emph{cerca}:\\
\emph{input}: matrice Soci e numero di tessera da cercare\\
\emph{output}: l’identificativo numerico associato al numero tessera
\end{frame}
